\documentclass[12pt]{article}
\usepackage[margin=1in, top=0.75in, bottom=0.75in]{geometry}
% worksheet-preamble.tex — the shipped preamble every generated document
% \input's. SOURCE OF TRUTH: this file. references/latex-templates.md documents
% these macros but no longer serves as the text agents transcribe — retyping a
% ~130-line preamble per document, three documents per run, is exactly the
% transcription-drift error class the rest of the pipeline exists to catch.
%
% Usage (compile.sh stages this file beside your .tex automatically):
%   \documentclass[12pt]{article}
%   \usepackage[margin=1in, top=0.75in, bottom=0.75in]{geometry}
%   % worksheet-preamble.tex — the shipped preamble every generated document
% \input's. SOURCE OF TRUTH: this file. references/latex-templates.md documents
% these macros but no longer serves as the text agents transcribe — retyping a
% ~130-line preamble per document, three documents per run, is exactly the
% transcription-drift error class the rest of the pipeline exists to catch.
%
% Usage (compile.sh stages this file beside your .tex automatically):
%   \documentclass[12pt]{article}
%   \usepackage[margin=1in, top=0.75in, bottom=0.75in]{geometry}
%   % worksheet-preamble.tex — the shipped preamble every generated document
% \input's. SOURCE OF TRUTH: this file. references/latex-templates.md documents
% these macros but no longer serves as the text agents transcribe — retyping a
% ~130-line preamble per document, three documents per run, is exactly the
% transcription-drift error class the rest of the pipeline exists to catch.
%
% Usage (compile.sh stages this file beside your .tex automatically):
%   \documentclass[12pt]{article}
%   \usepackage[margin=1in, top=0.75in, bottom=0.75in]{geometry}
%   \input{worksheet-preamble}
%   \input{figure-macros}          % when using the shipped figure macros
%
% geometry stays in the document, NOT here: worksheets use margin=1in
% (top/bottom 0.75in) while study guides use margin=0.85in (top/bottom 0.7in),
% and geometry options cannot be re-issued once the package is loaded.
%
% ASCII only in this file: pdflatex (the fallback engine) cannot typeset
% literal Unicode symbols, so documents must compile identically under both
% tectonic and pdflatex.

% ── Required packages (superset of the worksheet and study-guide shells; a
% worksheet loading mdframed it does not use costs nothing, one shared file
% that cannot drift is worth it) ─────────────────────────────────────────────
\usepackage{amsmath, amssymb}
\usepackage{tikz, pgfplots}
\usepackage{enumitem, fancyhdr, multicol, array, booktabs}
\usepackage{mdframed, xcolor}
\pgfplotsset{compat=1.18}
\usetikzlibrary{calc, angles, quotes}   % needed by the geometric figure macros

% ── Problem numbering ───────────────────────────────────────────────────────
\newcounter{prob}
\newcommand{\problem}[1]{\stepcounter{prob}\vspace{0.4cm}\noindent\textbf{\large\theprob.}\quad #1\vspace{0.3cm}}

% ── Shrink-to-fit title: \LARGE when it fits on one line, otherwise scale
% down to the text width — never enlarges, never wraps mid-phrase. (\resizebox
% comes from graphicx, which tikz already loads.) ────────────────────────────
\newlength{\fittedw}
\newcommand{\fittedtitle}[1]{%
  \settowidth{\fittedw}{\LARGE\bfseries #1}%
  \ifdim\fittedw>\linewidth\resizebox{\linewidth}{!}{\bfseries #1}%
  \else{\LARGE\bfseries #1}\fi}

% ── Running headers. Keep the left title SHORT (<= ~28 chars, e.g. "Triangle
% Trig Practice", not the full course name) — it shares the header line with
% the Name/Date blanks and a long title overlaps them. ──────────────────────
% \wsheader{Short Title} — student worksheet
\newcommand{\wsheader}[1]{%
  \pagestyle{fancy}\fancyhf{}%
  \fancyhead[L]{\textbf{\small #1}}%
  \fancyhead[R]{\small Name: \underline{\hspace{4.5cm}}~Date: \underline{\hspace{1.8cm}}}%
  \fancyfoot[C]{\thepage}%
  \renewcommand{\headrulewidth}{0.4pt}}

% \akheader{Short Topic} — answer key
\newcommand{\akheader}[1]{%
  \pagestyle{fancy}\fancyhf{}%
  \fancyhead[L]{\textbf{\small #1 --- Answer Key}}%
  \fancyhead[R]{\textit{\small For instructor/parent use}}%
  \fancyfoot[C]{\thepage}%
  \renewcommand{\headrulewidth}{0.4pt}}

% \ssheader{Topic} — skills summary / study guide
\newcommand{\ssheader}[1]{%
  \pagestyle{fancy}\fancyhf{}%
  \fancyhead[L]{\textbf{Skills Summary: #1}}%
  \fancyhead[R]{\small\textit{Study Guide \& Reference}}%
  \fancyfoot[C]{\thepage}%
  \renewcommand{\headrulewidth}{0.4pt}}

% ── Title blocks ────────────────────────────────────────────────────────────
% \wstitleblock{Title}{Course}{Date}
\newcommand{\wstitleblock}[3]{%
  \begin{center}
    \fittedtitle{#1}\\[0.3cm]
    {\large #2 \quad $\bullet$ \quad #3}
  \end{center}
  \noindent\rule{\linewidth}{0.4pt}\vspace{0.2cm}}

% \aktitleblock{Topic}{Course}{Date} — "Answer Key" is a subtitle BENEATH the
% topic, never appended to it: a long topic would wrap mid-phrase.
\newcommand{\aktitleblock}[3]{%
  \begin{center}
    \fittedtitle{#1}\\[0.15cm]
    {\Large\bfseries Answer Key}\\[0.25cm]
    {\large #2 \quad $\bullet$ \quad #3}
  \end{center}}

% \sstitleblock{Topic}
\newcommand{\sstitleblock}[1]{%
  \begin{center}
    {\LARGE\textbf{Skills Summary}}\\[0.2cm]
    {\large\textbf{#1}}\\[0.1cm]
    {\small\textit{Use this reference while completing your worksheet or when studying.}}
  \end{center}
  \vspace{0.1cm}
  \noindent\rule{\linewidth}{1.5pt}
  \vspace{0.3cm}}

% ── Study-guide color palette and boxes ─────────────────────────────────────
\definecolor{skillblue}{RGB}{30,100,180}
\definecolor{skillbluebg}{RGB}{235,244,255}
\definecolor{warnorange}{RGB}{200,90,0}
\definecolor{warnbg}{RGB}{255,243,230}
\definecolor{exgreen}{RGB}{20,120,60}
\definecolor{exgreenbg}{RGB}{230,248,238}

% Formula/rule box
\newmdenv[
  backgroundcolor=skillbluebg,
  linecolor=skillblue, linewidth=1.5pt,
  innertopmargin=6pt, innerbottommargin=6pt,
  innerleftmargin=10pt, innerrightmargin=10pt,
  skipabove=6pt, skipbelow=4pt
]{formulabox}

% Mini example box
\newmdenv[
  backgroundcolor=exgreenbg,
  linecolor=exgreen, linewidth=1pt,
  innertopmargin=5pt, innerbottommargin=5pt,
  innerleftmargin=10pt, innerrightmargin=10pt,
  skipabove=4pt, skipbelow=4pt
]{examplebox}

% Watch-out box
\newmdenv[
  backgroundcolor=warnbg,
  linecolor=warnorange, linewidth=1pt,
  innertopmargin=5pt, innerbottommargin=5pt,
  innerleftmargin=10pt, innerrightmargin=10pt,
  skipabove=4pt, skipbelow=4pt
]{watchoutbox}

% Skill section heading
\newcommand{\skillheading}[1]{%
  \vspace{0.4cm}
  {\large\textbf{\textcolor{skillblue}{#1}}}
  \vspace{0.1cm}
  \hrule height 1pt
  \vspace{0.2cm}
}

%   % figure-macros.tex — shipped geometric figure macros. SOURCE OF TRUTH: this
% file; references/latex-templates.md documents usage. Requires the tikz
% libraries loaded by worksheet-preamble.tex (calc, angles, quotes), so
% \input worksheet-preamble first.
%
% CONTRACT with the layout/prose checkers (do not break it):
%   - Macro names contain rt/tri/fig — that is how tests/check_layout.py and
%     tests/check_prose_consistency.py recognize a macro-style figure.
%   - Mandatory arguments are FLAT braces only (no nested {} inside an arg);
%     the checkers' arg scanners match [^{}]* per group.
%   - Value-bearing args carry exactly the printed label text; scale/styling
%     goes in the single optional [..] argument, which the checkers ignore
%     (so scale=1.2 is never mistaken for a figure value).
%
% Figure conventions (from SKILL.md / latex-templates.md): every number
% printed in a figure must come from the problem's verify JSON; vertices are
% A/B/C with sides a/b/c OPPOSITE them — the same convention verify.py's
% triangle type uses. Right triangles put the right angle at C, so c (= AB)
% is the hypotenuse — the SAME convention scripts/render_figures.py constructs
% (its RIGHT_ANGLE_VERTEX constant): \refrt sits beside renderer-built
% \probfig figures on the page, so the two sources must never disagree.
% tests/test_figure_convention.py enforces this by parsing the
% "% right angle mark" lines below — that trailing comment is load-bearing;
% keep it on every right-angle-mark \draw.
%
% ASCII only: pdflatex (the fallback engine) cannot typeset literal Unicode.

% ── \rtfig[scale]{bottom-leg label}{right-leg label}{hypotenuse label} ──────
% Right triangle, right angle at C (hypotenuse c = AB — the renderer's
% convention). A=(0,0), C=(4,0), B=(4,3): the 3-4-5 shape keeps every label
% clear of the drawn lines at default scale. Bottom leg = b (AC), right leg
% = a (CB), hypotenuse = c (AB) — name sides accordingly in the labels.
% Example: \rtfig{$8$}{$6$}{$x$}   or   \rtfig[0.9]{$b = 8$}{$a = 6$}{$c$}
\newcommand{\rtfig}[4][1.2]{%
  \begin{center}
  \begin{tikzpicture}[scale=#1]
    \coordinate (A) at (0,0);
    \coordinate (B) at (4,3);
    \coordinate (C) at (4,0);
    \draw[thick] (A) -- (C) -- (B) -- cycle;
    \draw (C) ++(-.25,0) -- ++(0,.25) -- ++(.25,0);  % right angle mark
    \node[below left] at (A) {$A$};
    \node[above right] at (B) {$B$};
    \node[below right] at (C) {$C$};
    \node[below] at (2,0) {#2};
    \node[right] at (4,1.5) {#3};
    \node[above left] at (2,1.5) {#4};
  \end{tikzpicture}
  \end{center}}

% ── \trifig[scale]{c}{b}{A-deg}{c-label}{b-label}{a-label}{angle-label} ─────
% General triangle, TO SCALE by construction (SAS): A at the origin, B on the
% x-axis at distance c, C computed from side b and angle A (TikZ trig takes
% degrees). Numeric args #2-#4 are the ACTUAL given values, so the drawing
% cannot disagree with the math; label args #5-#8 carry what is printed.
% Example: \trifig{7}{5}{34}{$c = 7$}{$b = 5$}{$a = ?$}{$34^\circ$}
\newcommand{\trifig}[8][0.6]{%
  \begin{center}
  \begin{tikzpicture}[scale=#1]
    \coordinate (A) at (0,0);
    \coordinate (B) at (#2,0);
    \coordinate (C) at ({#3*cos(#4)},{#3*sin(#4)});
    \draw[thick] (A) -- (B) -- (C) -- cycle;
    \node[below left]  at (A) {$A$};
    \node[below right] at (B) {$B$};
    \node[above]       at (C) {$C$};
    \node[below]       at ($(A)!0.5!(B)$) {#5};
    \node[above left]  at ($(A)!0.5!(C)$) {#6};
    \node[above right] at ($(B)!0.5!(C)$) {#7};
    \pic [draw, angle radius=7mm, angle eccentricity=1.6, "#8"] {angle = B--A--C};
  \end{tikzpicture}
  \end{center}}

% ── \refrt — the value-free reference figure ────────────────────────────────
% SKILL.md's figure-scope rule requires shared labelling conventions to live
% in ONE value-free reference figure placed with the directions; until now no
% template for it existed, so every generation invented one — and an invented
% figure that accidentally carried a number would itself trip the
% all-or-nothing figure rule it exists to satisfy. This macro has zero
% numerals in its output by construction (it takes no arguments and prints
% only letters), and the caption is baked in verbatim so the figure cannot be
% mistaken for a problem's givens. Right angle at C, so sides read: a = BC
% (opposite A), b = AC (opposite B), c = AB (opposite C, the hypotenuse) —
% matching render_figures.py's RIGHT_ANGLE_VERTEX, because this reference
% figure defines the convention students apply to every renderer-built
% figure on the sheet.
\newcommand{\refrt}{%
  \begin{center}
  \begin{tikzpicture}[scale=1.0]
    \coordinate (A) at (0,0);
    \coordinate (B) at (3.6,2.7);
    \coordinate (C) at (3.6,0);
    \draw[thick] (A) -- (C) -- (B) -- cycle;
    \draw (C) ++(-.25,0) -- ++(0,.25) -- ++(.25,0);  % right angle mark
    \node[below left] at (A) {$A$};
    \node[above right] at (B) {$B$};
    \node[below right] at (C) {$C$};
    \node[below] at (1.8,0) {$b$};
    \node[right] at (3.6,1.35) {$a$};
    \node[above left] at (1.8,1.35) {$c$};
    \pic [draw, angle radius=6mm, angle eccentricity=1.7, "$A$"] {angle = C--A--B};
  \end{tikzpicture}\\
  {\small\textit{How every triangle here is labelled. No values shown: use the
  numbers given in each problem.}}
  \end{center}}
          % when using the shipped figure macros
%
% geometry stays in the document, NOT here: worksheets use margin=1in
% (top/bottom 0.75in) while study guides use margin=0.85in (top/bottom 0.7in),
% and geometry options cannot be re-issued once the package is loaded.
%
% ASCII only in this file: pdflatex (the fallback engine) cannot typeset
% literal Unicode symbols, so documents must compile identically under both
% tectonic and pdflatex.

% ── Required packages (superset of the worksheet and study-guide shells; a
% worksheet loading mdframed it does not use costs nothing, one shared file
% that cannot drift is worth it) ─────────────────────────────────────────────
\usepackage{amsmath, amssymb}
\usepackage{tikz, pgfplots}
\usepackage{enumitem, fancyhdr, multicol, array, booktabs}
\usepackage{mdframed, xcolor}
\pgfplotsset{compat=1.18}
\usetikzlibrary{calc, angles, quotes}   % needed by the geometric figure macros

% ── Problem numbering ───────────────────────────────────────────────────────
\newcounter{prob}
\newcommand{\problem}[1]{\stepcounter{prob}\vspace{0.4cm}\noindent\textbf{\large\theprob.}\quad #1\vspace{0.3cm}}

% ── Shrink-to-fit title: \LARGE when it fits on one line, otherwise scale
% down to the text width — never enlarges, never wraps mid-phrase. (\resizebox
% comes from graphicx, which tikz already loads.) ────────────────────────────
\newlength{\fittedw}
\newcommand{\fittedtitle}[1]{%
  \settowidth{\fittedw}{\LARGE\bfseries #1}%
  \ifdim\fittedw>\linewidth\resizebox{\linewidth}{!}{\bfseries #1}%
  \else{\LARGE\bfseries #1}\fi}

% ── Running headers. Keep the left title SHORT (<= ~28 chars, e.g. "Triangle
% Trig Practice", not the full course name) — it shares the header line with
% the Name/Date blanks and a long title overlaps them. ──────────────────────
% \wsheader{Short Title} — student worksheet
\newcommand{\wsheader}[1]{%
  \pagestyle{fancy}\fancyhf{}%
  \fancyhead[L]{\textbf{\small #1}}%
  \fancyhead[R]{\small Name: \underline{\hspace{4.5cm}}~Date: \underline{\hspace{1.8cm}}}%
  \fancyfoot[C]{\thepage}%
  \renewcommand{\headrulewidth}{0.4pt}}

% \akheader{Short Topic} — answer key
\newcommand{\akheader}[1]{%
  \pagestyle{fancy}\fancyhf{}%
  \fancyhead[L]{\textbf{\small #1 --- Answer Key}}%
  \fancyhead[R]{\textit{\small For instructor/parent use}}%
  \fancyfoot[C]{\thepage}%
  \renewcommand{\headrulewidth}{0.4pt}}

% \ssheader{Topic} — skills summary / study guide
\newcommand{\ssheader}[1]{%
  \pagestyle{fancy}\fancyhf{}%
  \fancyhead[L]{\textbf{Skills Summary: #1}}%
  \fancyhead[R]{\small\textit{Study Guide \& Reference}}%
  \fancyfoot[C]{\thepage}%
  \renewcommand{\headrulewidth}{0.4pt}}

% ── Title blocks ────────────────────────────────────────────────────────────
% \wstitleblock{Title}{Course}{Date}
\newcommand{\wstitleblock}[3]{%
  \begin{center}
    \fittedtitle{#1}\\[0.3cm]
    {\large #2 \quad $\bullet$ \quad #3}
  \end{center}
  \noindent\rule{\linewidth}{0.4pt}\vspace{0.2cm}}

% \aktitleblock{Topic}{Course}{Date} — "Answer Key" is a subtitle BENEATH the
% topic, never appended to it: a long topic would wrap mid-phrase.
\newcommand{\aktitleblock}[3]{%
  \begin{center}
    \fittedtitle{#1}\\[0.15cm]
    {\Large\bfseries Answer Key}\\[0.25cm]
    {\large #2 \quad $\bullet$ \quad #3}
  \end{center}}

% \sstitleblock{Topic}
\newcommand{\sstitleblock}[1]{%
  \begin{center}
    {\LARGE\textbf{Skills Summary}}\\[0.2cm]
    {\large\textbf{#1}}\\[0.1cm]
    {\small\textit{Use this reference while completing your worksheet or when studying.}}
  \end{center}
  \vspace{0.1cm}
  \noindent\rule{\linewidth}{1.5pt}
  \vspace{0.3cm}}

% ── Study-guide color palette and boxes ─────────────────────────────────────
\definecolor{skillblue}{RGB}{30,100,180}
\definecolor{skillbluebg}{RGB}{235,244,255}
\definecolor{warnorange}{RGB}{200,90,0}
\definecolor{warnbg}{RGB}{255,243,230}
\definecolor{exgreen}{RGB}{20,120,60}
\definecolor{exgreenbg}{RGB}{230,248,238}

% Formula/rule box
\newmdenv[
  backgroundcolor=skillbluebg,
  linecolor=skillblue, linewidth=1.5pt,
  innertopmargin=6pt, innerbottommargin=6pt,
  innerleftmargin=10pt, innerrightmargin=10pt,
  skipabove=6pt, skipbelow=4pt
]{formulabox}

% Mini example box
\newmdenv[
  backgroundcolor=exgreenbg,
  linecolor=exgreen, linewidth=1pt,
  innertopmargin=5pt, innerbottommargin=5pt,
  innerleftmargin=10pt, innerrightmargin=10pt,
  skipabove=4pt, skipbelow=4pt
]{examplebox}

% Watch-out box
\newmdenv[
  backgroundcolor=warnbg,
  linecolor=warnorange, linewidth=1pt,
  innertopmargin=5pt, innerbottommargin=5pt,
  innerleftmargin=10pt, innerrightmargin=10pt,
  skipabove=4pt, skipbelow=4pt
]{watchoutbox}

% Skill section heading
\newcommand{\skillheading}[1]{%
  \vspace{0.4cm}
  {\large\textbf{\textcolor{skillblue}{#1}}}
  \vspace{0.1cm}
  \hrule height 1pt
  \vspace{0.2cm}
}

%   % figure-macros.tex — shipped geometric figure macros. SOURCE OF TRUTH: this
% file; references/latex-templates.md documents usage. Requires the tikz
% libraries loaded by worksheet-preamble.tex (calc, angles, quotes), so
% \input worksheet-preamble first.
%
% CONTRACT with the layout/prose checkers (do not break it):
%   - Macro names contain rt/tri/fig — that is how tests/check_layout.py and
%     tests/check_prose_consistency.py recognize a macro-style figure.
%   - Mandatory arguments are FLAT braces only (no nested {} inside an arg);
%     the checkers' arg scanners match [^{}]* per group.
%   - Value-bearing args carry exactly the printed label text; scale/styling
%     goes in the single optional [..] argument, which the checkers ignore
%     (so scale=1.2 is never mistaken for a figure value).
%
% Figure conventions (from SKILL.md / latex-templates.md): every number
% printed in a figure must come from the problem's verify JSON; vertices are
% A/B/C with sides a/b/c OPPOSITE them — the same convention verify.py's
% triangle type uses. Right triangles put the right angle at C, so c (= AB)
% is the hypotenuse — the SAME convention scripts/render_figures.py constructs
% (its RIGHT_ANGLE_VERTEX constant): \refrt sits beside renderer-built
% \probfig figures on the page, so the two sources must never disagree.
% tests/test_figure_convention.py enforces this by parsing the
% "% right angle mark" lines below — that trailing comment is load-bearing;
% keep it on every right-angle-mark \draw.
%
% ASCII only: pdflatex (the fallback engine) cannot typeset literal Unicode.

% ── \rtfig[scale]{bottom-leg label}{right-leg label}{hypotenuse label} ──────
% Right triangle, right angle at C (hypotenuse c = AB — the renderer's
% convention). A=(0,0), C=(4,0), B=(4,3): the 3-4-5 shape keeps every label
% clear of the drawn lines at default scale. Bottom leg = b (AC), right leg
% = a (CB), hypotenuse = c (AB) — name sides accordingly in the labels.
% Example: \rtfig{$8$}{$6$}{$x$}   or   \rtfig[0.9]{$b = 8$}{$a = 6$}{$c$}
\newcommand{\rtfig}[4][1.2]{%
  \begin{center}
  \begin{tikzpicture}[scale=#1]
    \coordinate (A) at (0,0);
    \coordinate (B) at (4,3);
    \coordinate (C) at (4,0);
    \draw[thick] (A) -- (C) -- (B) -- cycle;
    \draw (C) ++(-.25,0) -- ++(0,.25) -- ++(.25,0);  % right angle mark
    \node[below left] at (A) {$A$};
    \node[above right] at (B) {$B$};
    \node[below right] at (C) {$C$};
    \node[below] at (2,0) {#2};
    \node[right] at (4,1.5) {#3};
    \node[above left] at (2,1.5) {#4};
  \end{tikzpicture}
  \end{center}}

% ── \trifig[scale]{c}{b}{A-deg}{c-label}{b-label}{a-label}{angle-label} ─────
% General triangle, TO SCALE by construction (SAS): A at the origin, B on the
% x-axis at distance c, C computed from side b and angle A (TikZ trig takes
% degrees). Numeric args #2-#4 are the ACTUAL given values, so the drawing
% cannot disagree with the math; label args #5-#8 carry what is printed.
% Example: \trifig{7}{5}{34}{$c = 7$}{$b = 5$}{$a = ?$}{$34^\circ$}
\newcommand{\trifig}[8][0.6]{%
  \begin{center}
  \begin{tikzpicture}[scale=#1]
    \coordinate (A) at (0,0);
    \coordinate (B) at (#2,0);
    \coordinate (C) at ({#3*cos(#4)},{#3*sin(#4)});
    \draw[thick] (A) -- (B) -- (C) -- cycle;
    \node[below left]  at (A) {$A$};
    \node[below right] at (B) {$B$};
    \node[above]       at (C) {$C$};
    \node[below]       at ($(A)!0.5!(B)$) {#5};
    \node[above left]  at ($(A)!0.5!(C)$) {#6};
    \node[above right] at ($(B)!0.5!(C)$) {#7};
    \pic [draw, angle radius=7mm, angle eccentricity=1.6, "#8"] {angle = B--A--C};
  \end{tikzpicture}
  \end{center}}

% ── \refrt — the value-free reference figure ────────────────────────────────
% SKILL.md's figure-scope rule requires shared labelling conventions to live
% in ONE value-free reference figure placed with the directions; until now no
% template for it existed, so every generation invented one — and an invented
% figure that accidentally carried a number would itself trip the
% all-or-nothing figure rule it exists to satisfy. This macro has zero
% numerals in its output by construction (it takes no arguments and prints
% only letters), and the caption is baked in verbatim so the figure cannot be
% mistaken for a problem's givens. Right angle at C, so sides read: a = BC
% (opposite A), b = AC (opposite B), c = AB (opposite C, the hypotenuse) —
% matching render_figures.py's RIGHT_ANGLE_VERTEX, because this reference
% figure defines the convention students apply to every renderer-built
% figure on the sheet.
\newcommand{\refrt}{%
  \begin{center}
  \begin{tikzpicture}[scale=1.0]
    \coordinate (A) at (0,0);
    \coordinate (B) at (3.6,2.7);
    \coordinate (C) at (3.6,0);
    \draw[thick] (A) -- (C) -- (B) -- cycle;
    \draw (C) ++(-.25,0) -- ++(0,.25) -- ++(.25,0);  % right angle mark
    \node[below left] at (A) {$A$};
    \node[above right] at (B) {$B$};
    \node[below right] at (C) {$C$};
    \node[below] at (1.8,0) {$b$};
    \node[right] at (3.6,1.35) {$a$};
    \node[above left] at (1.8,1.35) {$c$};
    \pic [draw, angle radius=6mm, angle eccentricity=1.7, "$A$"] {angle = C--A--B};
  \end{tikzpicture}\\
  {\small\textit{How every triangle here is labelled. No values shown: use the
  numbers given in each problem.}}
  \end{center}}
          % when using the shipped figure macros
%
% geometry stays in the document, NOT here: worksheets use margin=1in
% (top/bottom 0.75in) while study guides use margin=0.85in (top/bottom 0.7in),
% and geometry options cannot be re-issued once the package is loaded.
%
% ASCII only in this file: pdflatex (the fallback engine) cannot typeset
% literal Unicode symbols, so documents must compile identically under both
% tectonic and pdflatex.

% ── Required packages (superset of the worksheet and study-guide shells; a
% worksheet loading mdframed it does not use costs nothing, one shared file
% that cannot drift is worth it) ─────────────────────────────────────────────
\usepackage{amsmath, amssymb}
\usepackage{tikz, pgfplots}
\usepackage{enumitem, fancyhdr, multicol, array, booktabs}
\usepackage{mdframed, xcolor}
\pgfplotsset{compat=1.18}
\usetikzlibrary{calc, angles, quotes}   % needed by the geometric figure macros

% ── Problem numbering ───────────────────────────────────────────────────────
\newcounter{prob}
\newcommand{\problem}[1]{\stepcounter{prob}\vspace{0.4cm}\noindent\textbf{\large\theprob.}\quad #1\vspace{0.3cm}}

% ── Shrink-to-fit title: \LARGE when it fits on one line, otherwise scale
% down to the text width — never enlarges, never wraps mid-phrase. (\resizebox
% comes from graphicx, which tikz already loads.) ────────────────────────────
\newlength{\fittedw}
\newcommand{\fittedtitle}[1]{%
  \settowidth{\fittedw}{\LARGE\bfseries #1}%
  \ifdim\fittedw>\linewidth\resizebox{\linewidth}{!}{\bfseries #1}%
  \else{\LARGE\bfseries #1}\fi}

% ── Running headers. Keep the left title SHORT (<= ~28 chars, e.g. "Triangle
% Trig Practice", not the full course name) — it shares the header line with
% the Name/Date blanks and a long title overlaps them. ──────────────────────
% \wsheader{Short Title} — student worksheet
\newcommand{\wsheader}[1]{%
  \pagestyle{fancy}\fancyhf{}%
  \fancyhead[L]{\textbf{\small #1}}%
  \fancyhead[R]{\small Name: \underline{\hspace{4.5cm}}~Date: \underline{\hspace{1.8cm}}}%
  \fancyfoot[C]{\thepage}%
  \renewcommand{\headrulewidth}{0.4pt}}

% \akheader{Short Topic} — answer key
\newcommand{\akheader}[1]{%
  \pagestyle{fancy}\fancyhf{}%
  \fancyhead[L]{\textbf{\small #1 --- Answer Key}}%
  \fancyhead[R]{\textit{\small For instructor/parent use}}%
  \fancyfoot[C]{\thepage}%
  \renewcommand{\headrulewidth}{0.4pt}}

% \ssheader{Topic} — skills summary / study guide
\newcommand{\ssheader}[1]{%
  \pagestyle{fancy}\fancyhf{}%
  \fancyhead[L]{\textbf{Skills Summary: #1}}%
  \fancyhead[R]{\small\textit{Study Guide \& Reference}}%
  \fancyfoot[C]{\thepage}%
  \renewcommand{\headrulewidth}{0.4pt}}

% ── Title blocks ────────────────────────────────────────────────────────────
% \wstitleblock{Title}{Course}{Date}
\newcommand{\wstitleblock}[3]{%
  \begin{center}
    \fittedtitle{#1}\\[0.3cm]
    {\large #2 \quad $\bullet$ \quad #3}
  \end{center}
  \noindent\rule{\linewidth}{0.4pt}\vspace{0.2cm}}

% \aktitleblock{Topic}{Course}{Date} — "Answer Key" is a subtitle BENEATH the
% topic, never appended to it: a long topic would wrap mid-phrase.
\newcommand{\aktitleblock}[3]{%
  \begin{center}
    \fittedtitle{#1}\\[0.15cm]
    {\Large\bfseries Answer Key}\\[0.25cm]
    {\large #2 \quad $\bullet$ \quad #3}
  \end{center}}

% \sstitleblock{Topic}
\newcommand{\sstitleblock}[1]{%
  \begin{center}
    {\LARGE\textbf{Skills Summary}}\\[0.2cm]
    {\large\textbf{#1}}\\[0.1cm]
    {\small\textit{Use this reference while completing your worksheet or when studying.}}
  \end{center}
  \vspace{0.1cm}
  \noindent\rule{\linewidth}{1.5pt}
  \vspace{0.3cm}}

% ── Study-guide color palette and boxes ─────────────────────────────────────
\definecolor{skillblue}{RGB}{30,100,180}
\definecolor{skillbluebg}{RGB}{235,244,255}
\definecolor{warnorange}{RGB}{200,90,0}
\definecolor{warnbg}{RGB}{255,243,230}
\definecolor{exgreen}{RGB}{20,120,60}
\definecolor{exgreenbg}{RGB}{230,248,238}

% Formula/rule box
\newmdenv[
  backgroundcolor=skillbluebg,
  linecolor=skillblue, linewidth=1.5pt,
  innertopmargin=6pt, innerbottommargin=6pt,
  innerleftmargin=10pt, innerrightmargin=10pt,
  skipabove=6pt, skipbelow=4pt
]{formulabox}

% Mini example box
\newmdenv[
  backgroundcolor=exgreenbg,
  linecolor=exgreen, linewidth=1pt,
  innertopmargin=5pt, innerbottommargin=5pt,
  innerleftmargin=10pt, innerrightmargin=10pt,
  skipabove=4pt, skipbelow=4pt
]{examplebox}

% Watch-out box
\newmdenv[
  backgroundcolor=warnbg,
  linecolor=warnorange, linewidth=1pt,
  innertopmargin=5pt, innerbottommargin=5pt,
  innerleftmargin=10pt, innerrightmargin=10pt,
  skipabove=4pt, skipbelow=4pt
]{watchoutbox}

% Skill section heading
\newcommand{\skillheading}[1]{%
  \vspace{0.4cm}
  {\large\textbf{\textcolor{skillblue}{#1}}}
  \vspace{0.1cm}
  \hrule height 1pt
  \vspace{0.2cm}
}

\wsheader{Centre and Spread}

\begin{document}
\wstitleblock{Measuring the Centre and Spread of a Data Set}{Grade 6--7}{}

\begin{center}\small\textit{Finding the mode, median and mean of a data set, measuring its
spread with the range and the interquartile range, and deciding which of those numbers
actually describes the group}\end{center}

\noindent\small\textbf{How to use this sheet.}
\begin{itemize}[label={$\bullet$}, leftmargin=1.1cm, itemsep=1pt, topsep=2pt]
  \item A \emph{measure of centre} answers ``what is a typical value here?'' --- mode, median
        and mean are three different answers to that one question.
  \item A \emph{measure of spread} answers ``how much do the values differ from each other?''
        --- the range and the interquartile range are two different answers to that one.
  \item Put the data in order before you touch the median, the quartiles or the range. On a
        dot plot that ordering is already done for you: read the dots from left to right.
  \item On a dot plot the label on the axis is a \emph{value} and each dot is one
        \emph{item}. Two dots above the $4$ means two items worth $4$ each.
\end{itemize}
\normalsize

\bigskip

\problem{\textbf{The mode.} The dot plot shows how many goals the Year 7 team scored in each of
its eleven games. Each dot is one game.
\begin{center}
\begin{tikzpicture}[x=0.8cm, y=0.33cm]
  \draw[thick] (-0.6,0) -- (5.6,0);
  \foreach \v in {0,...,5} {\draw (\v,0.4) -- (\v,-0.4);
     \node[below, font=\small] at (\v,-0.4) {\v};}
  \foreach \v/\n in {0/2, 1/3, 2/2, 3/2, 4/1, 5/1} {%
     \foreach \i in {1,...,\n} {\fill (\v,\i) circle (1.7pt);}}
  \node[font=\small] at (2.5,-3.4) {Goals scored per game};
\end{tikzpicture}
\end{center}
Which number of goals came up most often? That value is the \emph{mode}.
\par\vspace*{3.2cm}\ansline}

\problem{\textbf{The median.} Same eleven games.
\begin{center}
\begin{tikzpicture}[x=0.8cm, y=0.33cm]
  \draw[thick] (-0.6,0) -- (5.6,0);
  \foreach \v in {0,...,5} {\draw (\v,0.4) -- (\v,-0.4);
     \node[below, font=\small] at (\v,-0.4) {\v};}
  \foreach \v/\n in {0/2, 1/3, 2/2, 3/2, 4/1, 5/1} {%
     \foreach \i in {1,...,\n} {\fill (\v,\i) circle (1.7pt);}}
  \node[font=\small] at (2.5,-3.4) {Goals scored per game};
\end{tikzpicture}
\end{center}
Reading the dots from left to right already puts the eleven values in order, so the median
is the value the middle dot sits on. Count in from each end until you meet.
\par\smallskip
Which position is the middle one out of eleven?~\ansblank
\par\vspace*{3.4cm}\ansline}

\problem{\textbf{The mean.} Same eleven games.
\begin{center}
\begin{tikzpicture}[x=0.8cm, y=0.33cm]
  \draw[thick] (-0.6,0) -- (5.6,0);
  \foreach \v in {0,...,5} {\draw (\v,0.4) -- (\v,-0.4);
     \node[below, font=\small] at (\v,-0.4) {\v};}
  \foreach \v/\n in {0/2, 1/3, 2/2, 3/2, 4/1, 5/1} {%
     \foreach \i in {1,...,\n} {\fill (\v,\i) circle (1.7pt);}}
  \node[font=\small] at (2.5,-3.4) {Goals scored per game};
\end{tikzpicture}
\end{center}
Add the goals from all eleven games, then divide by the number of games. Two dots above the
$4$ would mean two games of four goals each --- count the dots, not the labels.
\par\smallskip
Total goals:~\ansblank \qquad Number of games:~\ansblank
\par\vspace*{4.4cm}\ansline}

\problem{\textbf{The range.} Same eleven games.
\begin{center}
\begin{tikzpicture}[x=0.8cm, y=0.33cm]
  \draw[thick] (-0.6,0) -- (5.6,0);
  \foreach \v in {0,...,5} {\draw (\v,0.4) -- (\v,-0.4);
     \node[below, font=\small] at (\v,-0.4) {\v};}
  \foreach \v/\n in {0/2, 1/3, 2/2, 3/2, 4/1, 5/1} {%
     \foreach \i in {1,...,\n} {\fill (\v,\i) circle (1.7pt);}}
  \node[font=\small] at (2.5,-3.4) {Goals scored per game};
\end{tikzpicture}
\end{center}
The range measures spread: how far apart the two most extreme games were.
\par\smallskip
Most goals:~\ansblank \qquad Fewest goals:~\ansblank
\par\vspace*{3.2cm}\ansline}

\problem{\textbf{Quartiles.} Same eleven games. Split the ordered list at the median, leaving
that middle value out of both halves.
\begin{center}
\begin{tikzpicture}[x=0.8cm, y=0.33cm]
  \draw[thick] (-0.6,0) -- (5.6,0);
  \foreach \v in {0,...,5} {\draw (\v,0.4) -- (\v,-0.4);
     \node[below, font=\small] at (\v,-0.4) {\v};}
  \foreach \v/\n in {0/2, 1/3, 2/2, 3/2, 4/1, 5/1} {%
     \foreach \i in {1,...,\n} {\fill (\v,\i) circle (1.7pt);}}
  \node[font=\small] at (2.5,-3.4) {Goals scored per game};
\end{tikzpicture}
\end{center}
The lower quartile $Q_1$ is the median of the lower half; the upper quartile $Q_3$ is the
median of the upper half. Write both halves out before you start.
\par\smallskip
$Q_1 =$~\ansblank \qquad $Q_3 =$~\ansblank
\par\vspace*{4.6cm}\noansline}

\problem{\textbf{The interquartile range.} Same eleven games.
\begin{center}
\begin{tikzpicture}[x=0.8cm, y=0.33cm]
  \draw[thick] (-0.6,0) -- (5.6,0);
  \foreach \v in {0,...,5} {\draw (\v,0.4) -- (\v,-0.4);
     \node[below, font=\small] at (\v,-0.4) {\v};}
  \foreach \v/\n in {0/2, 1/3, 2/2, 3/2, 4/1, 5/1} {%
     \foreach \i in {1,...,\n} {\fill (\v,\i) circle (1.7pt);}}
  \node[font=\small] at (2.5,-3.4) {Goals scored per game};
\end{tikzpicture}
\end{center}
The interquartile range is $Q_3 - Q_1$: the width of the middle half of the data.
\par\smallskip
\small Underneath your working, write one sentence saying what that number tells you about a
typical game.\normalsize
\par\vspace*{3.8cm}\ansline}

\problem{\textbf{One game at random.} Same eleven games. One of them is picked at random, and
every game is equally likely.
\begin{center}
\begin{tikzpicture}[x=0.8cm, y=0.33cm]
  \draw[thick] (-0.6,0) -- (5.6,0);
  \foreach \v in {0,...,5} {\draw (\v,0.4) -- (\v,-0.4);
     \node[below, font=\small] at (\v,-0.4) {\v};}
  \foreach \v/\n in {0/2, 1/3, 2/2, 3/2, 4/1, 5/1} {%
     \foreach \i in {1,...,\n} {\fill (\v,\i) circle (1.7pt);}}
  \node[font=\small] at (2.5,-3.4) {Goals scored per game};
\end{tikzpicture}
\end{center}
Count the \emph{games} that finished with $3$ or more goals --- not the labels on the axis.
Give the probability as a fraction.
\par\smallskip
Favourable games:~\ansblank \qquad Games in total:~\ansblank
\par\vspace*{3.6cm}\ansline}

\problem{\textbf{A different group.} Nine students recorded how many books they read over the
summer. Here are their totals, already in order, with the same data as a dot plot underneath.
\begin{center}\small
$3$, \quad $4$, \quad $5$, \quad $5$, \quad $6$, \quad $7$, \quad $8$, \quad $10$, \quad $24$
\normalsize\end{center}
\begin{center}
\begin{tikzpicture}[x=0.23cm, y=0.33cm]
  \draw[thick] (-1,0) -- (26,0);
  \foreach \v in {0,...,25} {\draw[gray] (\v,0.25) -- (\v,-0.25);}
  \foreach \v in {0,5,10,15,20,25} {\draw[thick] (\v,0.5) -- (\v,-0.5);
     \node[below, font=\small] at (\v,-0.5) {\v};}
  \foreach \v/\n in {3/1, 4/1, 5/2, 6/1, 7/1, 8/1, 10/1, 24/1} {%
     \foreach \i in {1,...,\n} {\fill (\v,\i) circle (1.7pt);}}
  \node[font=\small] at (12.5,-3.4) {Books read over the summer};
\end{tikzpicture}
\end{center}
Find the median number of books.
\par\smallskip
\small Which position is the middle one out of nine?\ \underline{\hspace{1.2cm}}\normalsize
\par\vspace*{3.4cm}\ansline}

\problem{\textbf{The mean of the same nine.}
\begin{center}\small
$3$, \quad $4$, \quad $5$, \quad $5$, \quad $6$, \quad $7$, \quad $8$, \quad $10$, \quad $24$
\normalsize\end{center}
\begin{center}
\begin{tikzpicture}[x=0.23cm, y=0.33cm]
  \draw[thick] (-1,0) -- (26,0);
  \foreach \v in {0,...,25} {\draw[gray] (\v,0.25) -- (\v,-0.25);}
  \foreach \v in {0,5,10,15,20,25} {\draw[thick] (\v,0.5) -- (\v,-0.5);
     \node[below, font=\small] at (\v,-0.5) {\v};}
  \foreach \v/\n in {3/1, 4/1, 5/2, 6/1, 7/1, 8/1, 10/1, 24/1} {%
     \foreach \i in {1,...,\n} {\fill (\v,\i) circle (1.7pt);}}
  \node[font=\small] at (12.5,-3.4) {Books read over the summer};
\end{tikzpicture}
\end{center}
Add all nine totals, then divide by nine.
\par\smallskip
Sum of the nine totals:~\ansblank
\par\vspace*{4.6cm}\ansline}

\problem{\textbf{The range of the same nine.}
\begin{center}\small
$3$, \quad $4$, \quad $5$, \quad $5$, \quad $6$, \quad $7$, \quad $8$, \quad $10$, \quad $24$
\normalsize\end{center}
\begin{center}
\begin{tikzpicture}[x=0.23cm, y=0.33cm]
  \draw[thick] (-1,0) -- (26,0);
  \foreach \v in {0,...,25} {\draw[gray] (\v,0.25) -- (\v,-0.25);}
  \foreach \v in {0,5,10,15,20,25} {\draw[thick] (\v,0.5) -- (\v,-0.5);
     \node[below, font=\small] at (\v,-0.5) {\v};}
  \foreach \v/\n in {3/1, 4/1, 5/2, 6/1, 7/1, 8/1, 10/1, 24/1} {%
     \foreach \i in {1,...,\n} {\fill (\v,\i) circle (1.7pt);}}
  \node[font=\small] at (12.5,-3.4) {Books read over the summer};
\end{tikzpicture}
\end{center}
Largest:~\ansblank \qquad Smallest:~\ansblank
\par\smallskip
\small Which single student decides this answer?\ \underline{\hspace{2.4cm}}\normalsize
\par\vspace*{3.2cm}\ansline}

\problem{\textbf{Devi is the student who read $24$.} These nine students have a mean of $8$
books and a median of $6$ books --- you worked both out above.
\begin{center}\small
$3$, \quad $4$, \quad $5$, \quad $5$, \quad $6$, \quad $7$, \quad $8$, \quad $10$, \quad $24$
\normalsize\end{center}
\begin{center}
\begin{tikzpicture}[x=0.23cm, y=0.33cm]
  \draw[thick] (-1,0) -- (26,0);
  \foreach \v in {0,...,25} {\draw[gray] (\v,0.25) -- (\v,-0.25);}
  \foreach \v in {0,5,10,15,20,25} {\draw[thick] (\v,0.5) -- (\v,-0.5);
     \node[below, font=\small] at (\v,-0.5) {\v};}
  \foreach \v/\n in {3/1, 4/1, 5/2, 6/1, 7/1, 8/1, 10/1, 24/1} {%
     \foreach \i in {1,...,\n} {\fill (\v,\i) circle (1.7pt);}}
  \node[font=\small] at (12.5,-3.4) {Books read over the summer};
\end{tikzpicture}
\end{center}
\textbf{(a)} Work out the mean of the other eight students, leaving Devi out.
\par\smallskip
Sum without Devi:~\ansblank \qquad Mean of the eight:~\ansblank
\par\smallskip
\textbf{(b)} What happens to the median when Devi is left out?\ \underline{\hspace{2.6cm}}
\par\smallskip
\textbf{(c)} Which measure of centre --- mean or median --- describes this group of readers
better? Write two or three sentences, and say what Devi's total does to each one.
\par\smallskip
\noindent\hrulefill
\par\vspace*{0.8cm}
\noindent\hrulefill
\par\vspace*{0.8cm}
\noindent\hrulefill
\par\vspace*{1.6cm}\noansline}

\problem{\textbf{Two ways to measure spread.} The same nine students, and the range you found
above.
\begin{center}\small
$3$, \quad $4$, \quad $5$, \quad $5$, \quad $6$, \quad $7$, \quad $8$, \quad $10$, \quad $24$
\normalsize\end{center}
\begin{center}
\begin{tikzpicture}[x=0.23cm, y=0.33cm]
  \draw[thick] (-1,0) -- (26,0);
  \foreach \v in {0,...,25} {\draw[gray] (\v,0.25) -- (\v,-0.25);}
  \foreach \v in {0,5,10,15,20,25} {\draw[thick] (\v,0.5) -- (\v,-0.5);
     \node[below, font=\small] at (\v,-0.5) {\v};}
  \foreach \v/\n in {3/1, 4/1, 5/2, 6/1, 7/1, 8/1, 10/1, 24/1} {%
     \foreach \i in {1,...,\n} {\fill (\v,\i) circle (1.7pt);}}
  \node[font=\small] at (12.5,-3.4) {Books read over the summer};
\end{tikzpicture}
\end{center}
\textbf{(a)} Find $Q_1$ and $Q_3$ for these nine values, then the interquartile range.
\par\smallskip
$Q_1 =$~\ansblank \qquad $Q_3 =$~\ansblank \qquad IQR:~\ansblank
\par\smallskip
\textbf{(b)} Which is larger, the range or the IQR?\ \underline{\hspace{2cm}}
\par\smallskip
\textbf{(c)} Which of the two gives a fairer picture of how spread out this group is? Write
two or three sentences saying what each one actually measures.
\par\smallskip
\noindent\hrulefill
\par\vspace*{0.8cm}
\noindent\hrulefill
\par\vspace*{0.8cm}
\noindent\hrulefill
\par\vspace*{1.6cm}\noansline}

\end{document}
